\documentclass[11pt]{article}
\input{../../shared/project_style.tex}
\rhead{Basics 2 Textbook Draft}

\begin{document}
\DocTitle{Basics 2: Vectors, Fields, and Coordinate Systems}{I - Mathematical and computational foundations}{Geometry and notation for Cartesian, cylindrical, toroidal, and field-aligned descriptions}

\begin{overviewbox}
\textbf{Textbook draft, version 1.0.} Plasma models are field theories. Correctly representing them requires a careful distinction among coordinates, basis vectors, vector components, metrics, and physical vectors. This chapter is designed for a reader with no prior plasma-physics training. It distinguishes exact identities from physical approximations and connects the central equations to later simulation modules.
\end{overviewbox}

\ProjectTOC

\section{Learning objectives}
After completing this note, the reader should be able to:
\begin{itemize}[leftmargin=1.6em]
\item Distinguish scalars, vectors, tensors, coordinates, and components.
\item Use dot, cross, and triple products physically and algebraically.
\item Represent vectors in Cartesian and cylindrical coordinates.
\item Explain why curvilinear basis vectors vary with position.
\item Use scale factors and Jacobians to compute distances, areas, and volumes.
\item Interpret radial, poloidal, toroidal, parallel, and perpendicular directions.
\end{itemize}

\section{Prerequisites and reading strategy}
\begin{itemize}[leftmargin=1.6em]
\item Basics 1 or equivalent algebra and elementary trigonometry.
\end{itemize}
On a first reading, emphasize physical meaning and dimensional checks. On a second reading, reproduce the derivations. Attempt the computational laboratory only after the limiting cases are understood.

\section{Scalars, vectors, tensors, and fields}

A scalar such as density $\rho(\mathbf x,t)$ assigns one number to each point. A vector such as velocity $\mathbf u(\mathbf x,t)$ assigns a magnitude and direction. A tensor such as the pressure tensor $\mathsf P$ describes a direction-dependent linear relation.

The physical vector is independent of coordinates, but its components are not:
\begin{equation}
\mathbf B=B^1\mathbf e_1+B^2\mathbf e_2+B^3\mathbf e_3.
\end{equation}
Changing the basis changes the components while leaving the vector unchanged. This is why a data file must identify both component ordering and coordinate system.


\section{Dot products, cross products, and projections}

The dot product
\begin{equation}
\mathbf a\cdot\mathbf b=ab\cos\theta
\end{equation}
measures projection and defines orthogonality. The component of $\mathbf a$ along a unit vector $\mathbf b$ is $a_\parallel=\mathbf a\cdot\mathbf b$.

The cross product
\begin{equation}
\mathbf a\times\mathbf b=ab\sin\theta\,\mathbf n
\end{equation}
is normal to the plane of the two vectors. The Lorentz magnetic force is perpendicular to both velocity and magnetic field. The scalar triple product
\begin{equation}
\mathbf a\cdot(\mathbf b\times\mathbf c)
\end{equation}
is an oriented volume and appears in coordinate Jacobians.


\section{Cartesian coordinates}

Cartesian coordinates use fixed orthonormal basis vectors:
\begin{equation}
\mathbf r=x\mathbf e_x+y\mathbf e_y+z\mathbf e_z,
\qquad
d\ell^2=dx^2+dy^2+dz^2.
\end{equation}
Because the basis vectors do not vary with position, differentiation is comparatively simple. Cartesian coordinates are ideal for local wave derivations and code-verification cases, though they do not conform naturally to a torus.


\section{Cylindrical coordinates and rotating basis vectors}

Cylindrical coordinates $(R,\phi,Z)$ satisfy
\begin{equation}
x=R\cos\phi,\qquad y=R\sin\phi,\qquad z=Z.
\end{equation}
The basis vectors are
\begin{align}
\mathbf e_R&=(\cos\phi,\sin\phi,0),\\
\mathbf e_\phi&=(-\sin\phi,\cos\phi,0),\\
\mathbf e_Z&=(0,0,1).
\end{align}
They obey
\begin{equation}
\pdv{\mathbf e_R}{\phi}=\mathbf e_\phi,
\qquad
\pdv{\mathbf e_\phi}{\phi}=-\mathbf e_R.
\end{equation}
The line element is
\begin{equation}
d\ell^2=dR^2+R^2d\phi^2+dZ^2,
\end{equation}
so an angular differential represents physical distance $R\,d\phi$.


\section{General coordinates, metric factors, and Jacobians}

For orthogonal coordinates $(q^1,q^2,q^3)$ with scale factors $h_i$,
\begin{equation}
d\ell^2=\sum_i h_i^2(dq^i)^2,
\qquad
dV=h_1h_2h_3\,dq^1dq^2dq^3.
\end{equation}
In cylindrical coordinates $(h_R,h_\phi,h_Z)=(1,R,1)$, so $dV=R\,dR\,d\phi\,dZ$. Omitting the Jacobian factor changes integrals, conservation laws, and matrix symmetry.

In nonorthogonal coordinates, the metric tensor
\begin{equation}
g_{ij}=\pdv{\mathbf r}{q^i}\cdot\pdv{\mathbf r}{q^j}
\end{equation}
determines distances and inner products. Covariant and contravariant components are then distinct.


\section{Toroidal and magnetic coordinates}

A simple circular torus may be written
\begin{align}
R&=R_0+r\cos\theta,\\
Z&=r\sin\theta,
\end{align}
where $\theta$ is poloidal and $\phi$ toroidal. Real stellarator surfaces require Fourier representations in both angles.

Magnetic coordinates commonly use a surface label $\psi$ together with angles. A magnetic surface satisfies
\begin{equation}
\mathbf B\cdot\nabla\psi=0.
\end{equation}
The field-aligned unit vector is $\mathbf b=\mathbf B/B$, and any vector decomposes as
\begin{equation}
\mathbf a=(\mathbf a\cdot\mathbf b)\mathbf b+
\left[\mathbf a-(\mathbf a\cdot\mathbf b)\mathbf b\right].
\end{equation}
The terms parallel and perpendicular always refer to a specified local magnetic field.

\begin{notebox}
A component called ``radial'' may mean cylindrical $R$, minor radius $r$, or flux coordinate $\psi$. Never infer the meaning from the word alone.
\end{notebox}


\section{Computational laboratory}

Generate points on a circular torus, convert between $(r,\theta,\phi)$ and Cartesian coordinates, and verify basis orthogonality. Approximate the torus volume using the correct Jacobian and compare with $2\pi^2R_0a^2$. Then omit the Jacobian deliberately and quantify the error. Implement identical coordinate conventions in Python and MATLAB.


\section{Summary}
\begin{itemize}[leftmargin=1.6em]
\item Vectors are coordinate independent, but their components and bases are not.
\item Curvilinear basis vectors and metric factors vary with position.
\item Jacobians are essential to physical integration and numerical conservation.
\item Toroidal, poloidal, radial, parallel, and perpendicular directions are distinct.
\end{itemize}

\SymbolsAcronymsSection
\begin{symtable}
$\mathbf e_i$ & basis vector & Direction associated with coordinate $q^i$. \\
$h_i$ & scale factor & Converts coordinate differential to physical length. \\
$g_{ij}$ & metric tensor & Inner products of coordinate basis vectors. \\
$R,\phi,Z$ & cylindrical coordinates & Radius, azimuth, and vertical coordinate. \\
$\mathbf b$ & field unit vector & $\mathbf B/B$. \\
\end{symtable}

\section{Exercises}
\begin{enumerate}[leftmargin=1.8em]
\item Prove the cylindrical-basis derivative identities.
\item Compute the physical length associated with $d\phi$ at two different major radii.
\item Decompose a given vector into parallel and perpendicular components.
\item Derive the volume of a circular torus from a coordinate mapping.
\item Explain the difference between a coordinate basis and an orthonormal physical basis.
\end{enumerate}

\section{References}
\begin{thebibliography}{99}
\bibitem{ref1} G. B. Arfken, H. J. Weber, and F. E. Harris, \emph{Mathematical Methods for Physicists}, 7th ed., Academic Press (2013).
\bibitem{ref2} G. Strang, \emph{Linear Algebra and Its Applications}, 4th ed., Brooks/Cole (2006).
\bibitem{ref3} W. A. Strauss, \emph{Partial Differential Equations: An Introduction}, 2nd ed., Wiley (2007).
\bibitem{ref4} F. F. Chen, \emph{Introduction to Plasma Physics and Controlled Fusion}, 3rd ed., Springer (2016).
\bibitem{ref5} R. J. Goldston and P. H. Rutherford, \emph{Introduction to Plasma Physics}, CRC Press (1995).
\bibitem{ref6} P. M. Bellan, \emph{Fundamentals of Plasma Physics}, Cambridge University Press (2006).
\end{thebibliography}

\end{document}
